\section{Position relative to established work}
\label{sec:positioning}

\subsection{The closest economic models}

Interdependence is already central to several models of AI automation. \citet{gans2026}
study \emph{O-Ring Automation}, in which task qualities are complementary rather than separable.
Automating one task changes the value of quality in the others. This is an important antecedent:
an automation decision cannot be evaluated only by the task it removes.

\citet{demirer2026} model ordered steps, manual execution, augmentation, automation, and job
boundaries shaped by handoff costs. Their empirical analysis (Section~7) also finds that a step is
more likely to be executed by AI when its neighbors in the production sequence are. Combining
O*NET tasks, published exposure labels, the Anthropic Economic Index, and model-generated workflow
orderings, this is an observational anchor for local propagation, not a causal estimate for
enterprise gates. The present paper examines how evidence-bearing handoffs affect downstream human
work, retaining verification, exceptions, and authorization in the comparison.

\citet{ide2025} embed AI in a knowledge economy with hierarchical production and problem
solving. The distinction between autonomous AI and assistance to humans changes organizational
and distributional outcomes. Our analysis treats a particular governance process as the observation
unit rather than solving for labor-market equilibrium. Table~\ref{tab:positioning} locates the
difference.

\begin{table}[!htb]
\centering\small
\caption{Closest models and the empirical object developed here.}
\label{tab:positioning}
\begin{tabularx}{\linewidth}{@{}P{3.1cm}LL@{}}
\toprule
Work & Central mechanism & Focus of this paper\\
\midrule
Gans and Goldfarb & Complementarity among task qualities. & A labor-accounting framework for
execution and assurance at organizational review interfaces.\\
Demirer et al. & Task chaining, handoff costs, and endogenous job boundaries. & Evidence
provenance, exception effort, and authorization within a governed process.\\
Ide and Talam\`as & AI in knowledge-based production hierarchies. & A fixed-process workload and
cost comparison before any equilibrium response.\\
\bottomrule
\end{tabularx}
\end{table}

\subsection{Established organizational foundations}

The gate is not a new managerial construct: Section~\ref{sec:contract} builds on Stage-Gate and on
artifact-centered process models. These approaches explain why an evolving project dossier
is a useful unit for tracing work across departments.

Knowledge hierarchies provide the antecedent for the human residual. \citet{garicano2000} models
specialization in solving problems; \citet[Sections~1 and~3]{garicano2015} explain
how ordinary problems are handled near production while specialists deal with exceptions, subject
to the cost of communicating and supplying help. \citet[pp.~775--777]{bainbridge1983} identifies the
supervisory, skill-retention, and recovery difficulties created by automation. The selection result of Section~\ref{sec:selection} connects these mechanisms to a measurable share of baseline labor and separates that
selection from the effort required after deployment.

Agentic Business Process Management addresses how human and software agents act within process
structures. \citet{vu2025} examine its development and practitioner perspectives.
\citet{calvanese2026} describe process frames, autonomy, explanation, interaction, and adaptation.
Our contribution is a measurement framework for the economic consequences of those implementations:
a producer-by-format comparison, the complete residual labor boundary, and a quality-constrained
forecast of workload.

\subsection{What the economics of the firm adds}

\citet[pp.~389--394]{coase1937} explains the comparison between coordinating transactions through
markets and organizing them within a firm. Review, negotiation, and internal authorization belong
to that coordination problem. Lowering their cost can change the feasible scale and depth of
internal control. The direction of the firm's boundary response depends on external as well as
internal coordination costs. \citet{shahidi2025} analyze the external channel: agents can reduce
search, communication, and contracting costs while introducing market frictions. Internal gate
savings alone therefore do not determine the firm's boundary.

Task-based economics remains complementary: substitution can coexist with new tasks and higher
demand \citep{autor2015,acemoglu2019}. Occupational exposure estimates concern
potential task effects \citep{eloundou2024}; measured AI assistance can have heterogeneous
productivity effects even within one occupation \citep{brynjolfsson2025}. The gate perspective
adds the links between these tasks and the work required to trust and authorize their outputs.

These works ground the representation and the mechanisms, but do not establish the universal endpoint
proposed here. The distinctive claim is the absence of a permanent human gate class across the gates an enterprise places in its DGF, illustrated here on three generic routes, including authority enactment and the support work introduced by automation. The paper
turns that claim into a contract-level thesis, supplies the labor account that measures progress toward
it, and specifies a measurement boundary that cannot hide the last human operators.
